% VonMises_on_econometrics.tex
% "Ludwig von Mises on Probability, Measurement, and Econometrics:
%  Passages for a Note on Model Ambiguity"
% Thomas J. Sargent, 2026

\documentclass[12pt]{article}

\usepackage{amsmath,amssymb}
\usepackage{booktabs}
\usepackage[margin=1.2in]{geometry}
\usepackage{setspace}
\usepackage{microtype}
\usepackage{hyperref}
\usepackage{natbib}
\usepackage{parskip}
\usepackage{csquotes}

\onehalfspacing

\hypersetup{
  colorlinks = true,
  linkcolor  = blue,
  citecolor  = blue,
  urlcolor   = blue
}

\title{Ludwig von Mises on Probability, Measurement,\\
       and Econometrics:\\
       Passages for a Note on Model Ambiguity}
\author{Thomas J.\ Sargent}
\date{August 2026}

\begin{document}
\maketitle

\begin{abstract}
This note collects passages from Ludwig von Mises on probability, measurement,
and econometrics.  Mises distinguished class probability from case probability,
denied that constant relations exist among economic magnitudes, and argued that
statistical series are historical records rather than estimates of parameters.
Section~\ref{sec:mill} collects what John Stuart Mill wrote about the
calculation of chances, because Mises quotes Mill and because Mill's chapter
says more than the quotation suggests.
I assembled the passages for use in work with Lars Peter Hansen on model
ambiguity.  Quotations are verbatim from the editions listed in the
bibliography.  Section~\ref{sec:use} states what Mises claimed and what he did
not claim.
\end{abstract}

%======================================================================
\section{Class Probability and Case Probability}
\label{sec:probability}

Chapter VI of \emph{Human Action} divides probability in two
\citep[pp.~105--118]{Mises1998HumanAction}.

\begin{quote}
There are two entirely different instances of probability; we may call them
class probability (or frequency probability) and case probability (or the
specific understanding of the sciences of human action).  The field for the
application of the former is the field of the natural sciences, entirely ruled
by causality; the field for the application of the latter is the field of the
sciences of human action, entirely ruled by teleology.
\hfill \citep[p.~107]{Mises1998HumanAction}
\end{quote}

The two definitions follow.

\begin{quote}
Class probability means: We know or assume to know, with regard to the problem
concerned, everything about the behavior of a whole class of events or
phenomena; but about the actual singular events or phenomena we know nothing
but that they are elements of this class.
\hfill \citep[p.~107]{Mises1998HumanAction}
\end{quote}

\begin{quote}
Case probability means: We know, with regard to a particular event, some of the
factors which determine its outcome; but there are other determining factors
about which we know nothing.
\hfill \citep[p.~110]{Mises1998HumanAction}
\end{quote}

Mises describes the calculus of probability as a translation of what the
analyst already knows.

\begin{quote}
For this defective knowledge the calculus of probability provides a
presentation in symbols of the mathematical terminology.  It neither expands
nor deepens nor complements our knowledge.  It translates it into mathematical
language.  Its calculations repeat in algebraic formulas what we knew
beforehand.
\hfill \citep[p.~108]{Mises1998HumanAction}
\end{quote}

He states the restriction on numerical evaluation flatly.

\begin{quote}
Case probability is not open to any kind of numerical evaluation.
\hfill \citep[p.~113]{Mises1998HumanAction}
\end{quote}

Section 5 of the chapter examines four statements a person might have made
before the 1944 presidential election, among them ``I estimate Roosevelt's
chances as 9 to 1.''  Mises grants that the statement means something and
denies that it is a probability:

\begin{quote}
It is a metaphorical expression.
\hfill \citep[p.~114]{Mises1998HumanAction}
\end{quote}

Section 4 closes with a passage that bears directly on decision making under
ambiguity.

\begin{quote}
The necessity to adjust his actions to other people's actions makes him a
speculator for whom success and failure depend on his greater or lesser ability
to understand the future.  Every investment is a form of speculation.  There is
in the course of human events no stability and consequently no safety.
\hfill \citep[p.~113]{Mises1998HumanAction}
\end{quote}

The same section separates three ways of confronting the future.  The gambler
knows the frequency of outcomes in a class and nothing about his own case.  The
engineer ``knows only soluble problems and problems which cannot be solved
under the present state of knowledge.''  The actor who must anticipate other
actors is neither \citep[p.~112]{Mises1998HumanAction}.

Mises also quotes John Stuart Mill on the misuse of the calculus.

\begin{quote}
The history of every branch of knowledge records instances of the misapplication
of the calculus of probability which, as John Stuart Mill observed, made it
``the real opprobrium of mathematics.''
\hfill \citep[pp.~106--107]{Mises1998HumanAction}
\end{quote}

%======================================================================
\section{John Stuart Mill on the Calculus of Chances}
\label{sec:mill}

Mises attributes the phrase ``the real opprobrium of mathematics'' to Mill.
The phrase occurs in \emph{A System of Logic}, Book III, Chapter XVIII, ``Of
the Calculation of Chances'' \citep{Mill1882Logic}.  Mill's chapter says four
things that bear on our subject.

\subsection{What Mill wrote}

Mill's sentence assigns a cause, and the cause is the sentence before it.

\begin{quote}
It is obvious, too, that even when the probabilities are derived from
observation and experiment, a very slight improvement in the data, by better
observations, or by taking into fuller consideration the special circumstances
of the case, is of more use than the most elaborate application of the calculus
to probabilities founded on the data in their previous state of inferiority.
The neglect of this obvious reflection has given rise to misapplications of the
calculus of probabilities which have made it the real opprobrium of mathematics.
\hfill \citep[Bk.~III, ch.~XVIII, \S3]{Mill1882Logic}
\end{quote}

The misapplications Mill has in mind are the eighteenth and nineteenth century
attempts to compute the credibility of witnesses and the reliability of jury
verdicts.  His complaint about them is a complaint about averages.

\begin{quote}
[T]he employment of it for such a purpose implies a misapprehension of the use
of averages, which serve, indeed, to protect those whose interest is at stake,
against mistaking the general result of large masses of instances, but are of
extremely small value as grounds of expectation in any one individual instance,
unless the case be one of those in which the great majority of individual
instances do not differ much from the average.
\hfill \citep[Bk.~III, ch.~XVIII, \S3]{Mill1882Logic}
\end{quote}

That is the distinction between class probability and case probability, stated
a century before \citet{Mises1998HumanAction} stated it.  Mill locates the
difficulty where Mises locates it, in the gap between a frequency in a class
and a ground for expectation about one member of it.

\subsection{Mill changed his mind}

The first edition of 1843 attacked the Laplacean theory.  Mill withdrew the
attack.

\begin{quote}
This view of the subject was taken in the first edition of the present work; but
I have since become convinced that the theory of chances, as conceived by
Laplace and by mathematicians generally, has not the fundamental fallacy which I
had ascribed to it.
\hfill \citep[Bk.~III, ch.~XVIII, \S1]{Mill1882Logic}
\end{quote}

\citet{Skyrms2018} traces the change to Mill's correspondence with John
Herschel and places it between the first edition and the second edition of
1846, hence before Herschel's review of Quetelet \citep{Herschel1850}.  What
Mill put in place of the attack is an epistemic reading of probability.

\begin{quote}
We must remember that the probability of an event is not a quality of the event
itself, but a mere name for the degree of ground which we, or some one else,
have for expecting it.  The probability of an event to one person is a different
thing from the probability of the same event to another, or to the same person
after he has acquired additional evidence.\ldots\ Every event is in itself
certain, not probable; if we knew all, we should either know positively that it
will happen, or positively that it will not.  But its probability to us means
the degree of expectation of its occurrence, which we are warranted in
entertaining by our present evidence.
\hfill \citep[Bk.~III, ch.~XVIII, \S1]{Mill1882Logic}
\end{quote}

Mises quotes the phrase from a chapter whose author had abandoned the position
that the phrase is usually taken to support.

\subsection{Mill's urn}

Mill argues \S2 from an urn of unknown composition.  A box holds balls of three
colors, white, black, and red, mixed in proportions the chooser does not know.
Mill concludes that the chooser should lay two to one against white, and he
adds a case that sharpens the problem.

\begin{quote}
[W]e might, on the contrary, know by authentic information that the box
contained ninety-nine balls of one color, and only one of the other; still, if
we are not told which color has only one, and which has ninety-nine, the drawing
of a white and of a black ball will be equally probable to us.
\hfill \citep[Bk.~III, ch.~XVIII, \S2]{Mill1882Logic}
\end{quote}

\citet{Ellsberg1961} put a three-color urn of unknown composition to the
opposite use.  Mill treats indifference as settling the matter.  Ellsberg used
the same apparatus to show that people decline to treat it as settled.

\subsection{Mill on statistical laws in the social sciences}

Book VI, Chapter IX of the \emph{Logic} states the condition under which a
statistical regularity may be carried to a new case.

\begin{quote}
But those immediate causes depend on remote causes; and the empirical law,
obtained by this indirect mode of observation, can only be relied on as
applicable to unobserved cases, so long as there is reason to think that no
change has taken place in any of the remote causes on which the immediate causes
depend.  In making use, therefore, of even the best statistical generalizations
for the purpose of inferring (though it be only conjecturally) that the same
empirical laws will hold in any new case, it is necessary that we be well
acquainted with the remoter causes, in order that we may avoid applying the
empirical law to cases which differ in any of the circumstances on which the
truth of the law ultimately depends.  And thus, even where conclusions derived
from specific observation are available for practical inferences in new cases,
it is necessary that the deductive science should stand sentinel over the whole
process.
\hfill \citep[Bk.~VI, ch.~IX, \S5]{Mill1882Logic}
\end{quote}

The next section denies that enough data can be assembled to average over the
combinations that matter.

\begin{quote}
Now the number of instances necessary to exhaust the whole round of combinations
of the various influential circumstances, and thus afford a fair average, never
can be obtained.  Not only we can never learn with sufficient authenticity the
facts of so many instances, but the world itself does not afford them in
sufficient numbers, within the limits of the given state of society and
civilization which such inquiries always presuppose.
\hfill \citep[Bk.~VI, ch.~IX, \S6]{Mill1882Logic}
\end{quote}

\subsection{Mill's remedy}

Mill and Mises identify the same difficulty.  A fitted regularity holds only
while the causes behind it hold, and the analyst cannot verify that they hold.
Their remedies differ.  Mill asks for better observations and for a deductive
science that stands sentinel over the inference.  Mises concludes that
quantification in economics is impossible.  Mill's remedy is the one that a
decision maker who distrusts his model can act on.

%======================================================================
\section{Constant Relations and Measurement}
\label{sec:constants}

Chapter II, section 10 of \emph{Human Action} contains the argument on which
the rest depends.

\begin{quote}
There are, in the field of economics, no constant relations, and consequently
no measurement is possible.  If a statistician determines that a rise of 10 per
cent in the supply of potatoes in Atlantis at a definite time was followed by a
fall of 8 per cent in the price, he does not establish anything about what
happened or may happen with a change in the supply of potatoes in another
country or at another time.  He has not ``measured'' the ``elasticity of
demand'' of potatoes.  He has established a unique and individual historical
fact.
\hfill \citep[p.~55]{Mises1998HumanAction}
\end{quote}

\begin{quote}
The impracticability of measurement is not due to the lack of technical methods
for the establishment of measure.  It is due to the absence of constant
relations.  If it were only caused by technical insufficiency, at least an
approximate estimation would be possible in some cases.  But the main fact is
that there are no constant relations.  Economics is not, as ignorant positivists
repeat again and again, backward because it is not ``quantitative.''  It is not
quantitative and does not measure because there are no constants.  Statistical
figures referring to economic events are historical data.  They tell us what
happened in a nonrepeatable historical case.
\hfill \citep[p.~56]{Mises1998HumanAction}
\end{quote}

%======================================================================
\section{Econometrics Named}
\label{sec:econometrics}

\subsection{The motto of the Econometric Society}

Chapter XVI, section 5 of \emph{Human Action} treats three varieties of
mathematical economics.  The first is the statistical variety.

\begin{quote}
The first variety is represented by the statisticians who aim at discovering
economic laws from the study of economic experience.  They aim to transform
economics into a ``quantitative'' science.  Their program is condensed in the
motto of the Econometric Society: Science is measurement.
\hfill \citep[p.~347]{Mises1998HumanAction}
\end{quote}

\begin{quote}
Statistics is a method for the presentation of historical facts concerning
prices and other relevant data of human action.  It is not economics and cannot
produce economic theorems and theories.  The statistics of prices is economic
history.\ldots\ Nobody ever was or ever will be in a position to observe a
change in one of the market data ceteris paribus.  There is no such thing as
quantitative economics.  All economic quantities we know about are data of
economic history.
\hfill \citep[pp.~347--348]{Mises1998HumanAction}
\end{quote}

Mises then takes up \citet{Schultz1938} on the elasticity of demand and
concludes that Schultz's results ``do not refer to potatoes in general, but to
potatoes in the United States in the years from 1875 to 1929''
\citep[p.~349]{Mises1998HumanAction}.  Later in the same chapter he contrasts
economics with mechanics.

\begin{quote}
Our modern industrial civilization is mainly an accomplishment of this
utilization of the differential equations of physics.  No such constant
relations exist, however, between economic elements.  The equations formulated
by mathematical economics remain a useless piece of mental gymnastics and would
remain so even if they were to express much more than they really do.
\hfill \citep[p.~351]{Mises1998HumanAction}
\end{quote}

\subsection{Mises quotes the Cowles Commission}

The Introduction to \emph{Theory and History} contains the passage closest to
our subject.  Mises quotes the Cowles Commission's own account of its
assumptions.

\begin{quote}
The econometrician is unable to disprove this fact, which cuts the ground from
under his reasoning.  He cannot help admitting that there are no ``behavior
constants.''  Nonetheless he wants to introduce some numbers, arbitrarily chosen
on the basis of a historical fact, as ``unknown behavior constants.''  The sole
excuse he advances is that his hypotheses are ``saying only that these unknown
numbers remain reasonably constant through a period of years.''
\hfill \citep[p.~11]{Mises1957TheoryHistory}
\end{quote}

The internal quotation is from \citet{Cowles1949}, page 7.  Mises continues:

\begin{quote}
Now whether such a period of supposed constancy of a definite number is still
lasting or whether a change in the number has already occurred can only be
established later on.\ldots\ It is the assertion of a historical fact, not of a
constant that can be resorted to in attempts to predict future events.\ldots\
Nonhuman entities react according to regular patterns; man chooses.
\hfill \citep[p.~11]{Mises1957TheoryHistory}
\end{quote}

This paragraph states the parameter-instability half of the argument later made
in \citet{Lucas1976}.  Mises directs it at the Cowles program and sources it to
a Cowles document.

\subsection{The 1962 formulation}

\emph{The Ultimate Foundation of Economic Science} states the case at its
sharpest.

\begin{quote}
Deluded by the idea that the sciences of human action must ape the technique of
the natural sciences, hosts of authors are intent upon a quantification of
economics.\ldots\ Their motto is the positivistic maxim: Science is measurement.
Supported by rich funds, they are busy reprinting and rearranging statistical
data provided by governments, by trade associations, and by corporations and
other enterprises.  They try to compute the arithmetical relations among various
of these data and thus to determine what they call, by analogy with the natural
sciences, correlations and functions.  They fail to realize that in the field of
human action statistics is always history\ldots\ As a method of economic
analysis econometrics is a childish play with figures that does not contribute
anything to the elucidation of the problems of economic reality.
\hfill \citep[p.~63]{Mises1962Ultimate}
\end{quote}

%======================================================================
\section{Forecasting}
\label{sec:forecasting}

Chapter 4 of \emph{The Ultimate Foundation of Economic Science} treats
uncertainty directly.

\begin{quote}
Technological engineering does not eliminate the aleatory element of human
existence; it merely restricts its field a little.  There always remains an
orbit that to the limited knowledge of man appears as an orbit of pure chance
and marks life as a gamble.  Man and his works are always exposed to the impact
of unforeseen and uncontrollable events.
\hfill \citep[pp.~64--65]{Mises1962Ultimate}
\end{quote}

Mises notes that a published forecast changes the process it forecasts.

\begin{quote}
Forecasting may prove mistaken if it induces men to proceed successfully in a
way that is designed to avoid the happening of the forecast events.\ldots\ If
people are told in May that the boom going on will continue for several months
and will not end in a crash before December, they will try to sell as soon as
possible, at any rate before December.  Then the boom will come to an end before
the day indicated by the prediction.
\hfill \citep[p.~66]{Mises1962Ultimate}
\end{quote}

He grants economics a predictive role and bounds it.

\begin{quote}
Economics can predict the effects to be expected from resorting to definite
measures of economic policies.\ldots\ But, of course, this prediction can be
only ``qualitative.''  It cannot be ``quantitative'' as there are no constant
relations between the factors and effects concerned.
\hfill \citep[p.~67]{Mises1962Ultimate}
\end{quote}

%======================================================================
\section{Mises and the Statisticians}
\label{sec:statisticians}

Richard von Mises (1883--1953), Ludwig's younger brother, developed the
frequency theory of probability from 1919 onward \citep{MisesR1957Probability}.
Ludwig accepted that theory for class probability.  His claim is not that
frequency probability is wrong but that it says nothing about singular events.
The editorial apparatus of the Scholar's Edition argues that Chapter VI of
\emph{Human Action} was stimulated by Richard's work
\citep{Mises1998HumanAction}.

Gerhard Tintner and Oskar Morgenstern both attended the \emph{Privatseminar}
that Mises conducted at the Vienna Chamber of Commerce
\citep{Haberler1961,Hulsmann2007}.  Tintner became one of the founders of
econometrics as a field \citep{Tintner1952}.  Several sources describe Mises as
the doctoral supervisor of both men.  Mises held no salaried chair at Vienna;
he was \emph{Privatdozent} from 1913 and unsalaried \emph{ausserordentlicher
Professor} from 1918, and other accounts name Hans Mayer as Morgenstern's
principal academic supervisor.  The seminar attendance is documented.  The
supervision should be checked against \citet{Hulsmann2007} and the Vienna
faculty records before it is asserted in print.

%======================================================================
\section{What Mises Claimed}
\label{sec:use}

Mises made three claims that should be kept separate.

First, economic time series record singular historical events rather than
repeated draws from a stable law.  A fitted coefficient summarises a time and a
place.  Mises is right about this, and the point is not a minor one.  It is the
same observation that motivates \citet{Lucas1976} and that motivates our
treatment of a decision maker who does not know the data generating process
\citep{HansenSargent2008}.

Second, the analyst who writes down a constant is making an assumption about
the future that his data cannot support.  The Cowles passage in
\emph{Theory and History} is the strongest form of the argument, because there
Mises is not asserting an impossibility.  He is pointing at a specific
assumption that the practitioners themselves wrote down.

Third, quantification in economics is therefore impossible.  This claim is
stronger than the first two and does not follow from them.  A decision maker
who distrusts his model need not abandon numbers.  He can carry a set of models
and act well against the worst of them.  That is the program in
\citet{HansenSargent2008}, and it addresses the problem that Mises identified
rather than the conclusion he drew from it.

Mises worked before the tools for the third step existed.  Frank Knight's
distinction between measurable risk and unmeasurable uncertainty
\citep{Knight1921} was the state of the art when Mises wrote Chapter VI, and
Mises sharpened it by locating the difficulty in the singularity of the event
rather than in the analyst's ignorance.  He posed the problem well.  We should
say so, and we should say that the answer we give is one he did not have
available.

Mill reached the same difficulty and drew the opposite conclusion.  Quoting
Mill on the opprobrium of mathematics alongside Mises requires reporting that
Mill withdrew his attack on the Laplacean theory, and that Mill asked for
better data under the supervision of theory rather than for less measurement
\citep[Bk.~III, ch.~XVIII; Bk.~VI, ch.~IX]{Mill1882Logic}.

\bigskip
\bibliographystyle{plainnat}
\bibliography{VonMises_on_econometrics}

\end{document}
